\section{Implementarea use case-urilor de aplicație}
Use case-urile reprezintă centrul implementării. Ele primesc comenzi clare și returnează rezultate sau erori controlate. De exemplu, un use case pentru aprobarea unui candidat la recompensă nu trebuie să primească un request HTTP brut, ci actorul, identificatorul candidatului, decizia și motivul. Astfel, aceeași logică poate fi apelată dintr-o rută HTTP, dintr-un test de integrare sau dintr-un job administrativ. Această separare scade dependența de framework și crește capacitatea de testare.

Un model important este injectarea porturilor. Use case-ul nu cunoaște detalii despre Diesel sau Ethereum, ci lucrează cu interfețe care exprimă nevoi: găsirea candidatului, salvarea auditului, verificarea permisiunilor, crearea unei tranzacții sau trimiterea unei notificări. Implementările concrete sunt furnizate de stratul de infrastructură. Pentru o disertație, acest model arată maturitate arhitecturală: codul nu este doar funcțional, ci organizat pentru evoluție.

\begin{lstlisting}[style=cod,caption={Structura conceptuală a unui use case.},label={lst:usecase}]
pub async fn handle(command: ApproveRewardCommand) -> Result<RewardOutput, RewardError> {
    access_control.ensure_can_review(command.actor, command.course_id).await?;
    let candidate = rewards.find_candidate(command.candidate_id).await?;
    let decision = domain::approve_candidate(candidate, command.reason)?;
    rewards.save_decision(decision).await?;
    notifications.reward_decision(command.actor, command.candidate_id).await?;
    Ok(RewardOutput::approved(command.candidate_id))
}
\end{lstlisting}

Fragmentul este conceptual, dar descrie stilul implementării: verificare de acces, încărcare stare, regulă de domeniu, persistență, efect secundar controlat și răspuns. Această ordine reduce riscul ca sistemul să execute efecte externe înainte de a valida regulile principale. În fluxurile reale, unele operații sunt încapsulate în tranzacții, iar efectele externe sunt mutate în joburi sau etape ulterioare.

\section{Idempotentă și consistență}
Idempotenta este o proprietate critică pentru operațiile care pot fi repetate de utilizator, de browser, de un retry HTTP sau de un worker. Dacă un cursant apasă de două ori pe trimiterea evaluării, sistemul nu trebuie să creeze două recompense pentru același eveniment. Dacă un worker reia un job după eșec, nu trebuie să crediteze wallet-ul de două ori. Dacă un indexer vede același eveniment blockchain după restart, trebuie să îl recunoască drept deja procesat.

RustLearn gestionează idempotenta prin chei unice, statusuri și tranziții controlate. Cheia de idempotentă poate include actorul, sursa, evenimentul și resursa. Pentru recompense, acest lucru înseamnă că același eveniment educațional produce cel mult un candidat valid. Pentru depuneri de tokenuri, evenimentul on-chain este identificat prin chain id, contract, hash, index de log și adrese relevante. Astfel, sistemul poate relua procesarea fără a produce efecte duplicate.

Consistența trebuie tratată diferit pentru operațiile locale și cele externe. În PostgreSQL, o tranzacție poate proteja mai multe scrieri. În Ethereum, confirmarea este externă și poate întârzia. De aceea, aplicația nu presupune că totul se finalizează într-un singur pas. Ea creează stări intermediare și mecanisme de reconciliere. Această abordare este mai realistă decât încercarea de a transforma blockchain-ul într-o extensie atomică a bazei de date.

\section{Implementarea permisiunilor pe scopuri}
Permisiunile sunt implementate pe mai multe niveluri. La nivel de platformă, ele controlează operații precum aprobarea aplicațiilor de profesor, gestionarea politicilor, fraud blocks și exporturile globale. La nivel de organizație, ele controlează membrii, cursurile sponsorizate și rapoartele organizației. La nivel de curs, ele controlează autoratul, publicarea, înscrierile, studenții și candidații la recompense. Această structură corespunde situațiilor reale din educație, unde autoritatea este contextuală.

Middleware-urile verifică permisiunile înainte ca handler-ul să apeleze use case-ul sensibil. Totuși, verificarea nu se limitează la middleware. Use case-urile importante repetă sau specializează validarea atunci când decizia depinde de starea resursei. Această dublă atenție este justificată deoarece permisiunile sunt o zonă de risc major. Interfața poate ascunde un buton, middleware-ul poate bloca ruta, iar use case-ul poate refuza operația dacă resursa nu se află în starea corectă.

\section{Implementarea recompenselor ca proces}
Subdomeniul recompenselor este unul dintre cele mai elaborate deoarece combină pedagogie, economie și securitate. Implementarea pornește de la candidatul la recompensă. Candidatul conține sursa, evenimentul, actorii relevanți, starea și cheia de idempotentă. După creare, el poate fi evaluat de profesor, apoi de platformă. În această etapă se stabilește suma aprobată și politica aplicată. Fiecare decizie produce audit, astfel încât istoricul să poată fi reconstruit.

După aprobare, sistemul creează o intenție de plată sau un job de execuție. Această separare este importantă deoarece o plată tokenizată poate depinde de disponibilitatea contractului, de nonce, de gas, de semnătură și de confirmarea rețelei. Dacă etapa token reușește, portofelul intern este creditat. Dacă portofelul este creditat, utilizatorul primește notificare. Dacă una dintre etape eșuează, candidatul poate intra într-o stare de reconciliere, nu într-o stare ambiguă.

\begin{table}[H]
    \centering
    \caption{Etapele implementate pentru recompensa LearnToken (elaborare proprie).}
    \label{tab:etape_recompensa_impl}
    \begin{tabular}{P{3.3cm}P{5.2cm}P{5.2cm}}
        \toprule
        Etapă & Responsabilitate & Motiv arhitectural \\
        \midrule
        Candidat & Leagă evenimentul educațional de recompensă & Separă învățarea de plata efectivă \\
        Decizie profesor & Validează meritul educațional & Păstrează context pedagogic \\
        Aprobare sumă & Aplică politică și control platformă & Previne recompense necontrolate \\
        Execuție token & Confirmă operația cu LearnToken & Adaugă audit extern \\
        Creditare wallet & Reflectă valoarea în aplicație & Permite dashboard și rapoarte \\
        Reconciliere & Repară efecte parțiale & Evită dublarea și pierderea stării \\
        \bottomrule
    \end{tabular}
\end{table}

\section{Integrarea contractelor inteligente}
Contractele sunt integrate prin două direcții: deploy și apeluri operaționale. La pornire sau în mediul local, aplicația poate inițializa adresele necesare pentru LearnToken și contractele auxiliare. Pentru operațiile de import sau transfer semnat, backend-ul pregătește datele, respectă chain id-ul, folosește ABI-ul generat și verifică rezultatele. Contractele rămân responsabile pentru reguli on-chain precum owner, nonce, permit și validarea semnăturii.

O decizie importantă este folosirea standardelor existente. ERC-20 oferă interoperabilitate, OpenZeppelin reduce riscul implementării greșite a funcțiilor de bază, iar EIP-712/EIP-2612 îmbunătățesc experiența semnăturilor \cite{openzeppelinERC20,eip712,eip2612}. Sistemul nu inventează un token complet nou, ci extinde un model cunoscut. Această alegere este defensivă și justificată academic: într-un domeniu cu risc financiar, reutilizarea standardelor mature este preferabilă originalității inutile.

\section{Implementarea portofelelor}
Portofelul intern este o oglindă operațională a valorii atribuite utilizatorului sau organizației în platformă. El nu înlocuiește portofelul blockchain, ci îl completează. Utilizatorul poate avea o adresă externă, dar aplicația are nevoie și de un ledger intern pentru istoricul recompenselor, taxelor, compensărilor, depunerilor și retragerilor. Acest ledger permite rapoarte, filtrare și reconciliere rapidă.

Implementarea portofelelor diferențiază actorii. Un wallet de utilizator are alt context decât un wallet de organizație. Organizațiile pot sponsoriza programe și pot avea audit financiar separat. De asemenea, platforma poate aplica taxe pentru anumite operații dacă suportă costul de gas. Această modelare arată că portofelul nu este doar o adresă Ethereum, ci o componentă a produsului educațional.

\section{Worker-ul și toleranța la eșec}
Worker-ul este construit pentru operații care nu trebuie să blocheze răspunsul HTTP. Procesarea video este exemplul principal: fișierul poate fi mare, transcodarea poate dura, iar rezultatul trebuie încărcat în S3. Dacă această operație ar rula în cererea profesorului, platforma ar fi fragilă. Prin joburi asincrone, profesorul primește o stare, iar sistemul poate procesa materialul în fundal.

Toleranța la eșec se bazează pe statusuri, retry-uri și heartbeat. Un job poate eșua temporar din cauza unui fișier lipsă, a unei probleme S3 sau a unei erori de transcodare. Sistemul trebuie să păstreze eroarea și să permită reluarea controlată. Heartbeat-ul arată că worker-ul este activ, iar readiness-ul general poate indica dacă platforma are toate dependențele necesare. Astfel, operarea nu depinde doar de loguri manuale.

\section{Implementarea frontend-ului pe stări}
Frontend-ul nu trebuie să presupună că toate datele sunt disponibile instantaneu. Fiecare ecran are stări de încărcare, eroare, conținut gol și conținut disponibil. Acest model este important pentru interfețele educaționale deoarece utilizatorul trebuie să înțeleagă de ce nu vede un curs, de ce nu poate aproba o recompensă sau de ce un wallet nu afișează încă tranzacția. O stare goală bine formulată poate preveni confuzia mai eficient decât o simplă listă vidă.

Interfața citește sesiunea curentă și capabilitățile asociate. Acțiunile sunt afișate în funcție de permisiuni, nu doar de tipul paginii. De exemplu, un utilizator aflat în zona \texttt{/teach} poate avea acces la un curs, dar nu neapărat la toate operațiile de recompensare. Aceeași logică se aplică pentru organizații și admin. Astfel, frontend-ul devine o extensie a modelului de permisiuni, nu un strat independent care poate crea așteptări false.

\section{KYC, audit și operații administrative}
KYC-ul este inclus ca flux administrativ deoarece recompensele tokenizate pot cere verificări suplimentare în scenarii reale. Chiar dacă implementarea de disertație poate fi locală, arhitectura prevede statusuri, review, audit și decizii. Acest lucru permite extinderea ulterioară către furnizori reali fără a schimba toate fluxurile de portofel.

Auditul este prezent în mai multe domenii: aplicații de profesor, organizații, recompense, politici și portofele. El nu este doar o listă de loguri tehnice, ci o parte a modelului de produs. Când o decizie este contestată, platforma trebuie să poată răspunde cine a decis, când, cu ce motiv și ce stare avea resursa. Această capacitate este esențială pentru încrederea într-un sistem care combină educația cu tokenuri.

\section{Observabilitate și administrare locală}
Implementarea include endpoint-uri de health și readiness, comenzi Make și mediu Docker Compose. Aceste elemente creează o experiență reproductibilă pentru dezvoltare și verificare. \texttt{make test}, \texttt{make preflight}, \texttt{make dev}, \texttt{make dev-worker} și comenzile de migrare definesc pași clari. Într-un proiect aplicativ, această claritate este aproape la fel de importantă ca implementarea funcțiilor, deoarece permite evaluatorului sau dezvoltatorului să reproducă rezultatul.

Observabilitatea poate fi extinsă în viitor cu metrici, tracing și dashboard-uri. Totuși, chiar și forma actuală arată o direcție corectă: aplicația nu pornește ca un proces opac, ci își declară starea dependențelor. Worker-ul nu este presupus sănătos doar pentru că procesul există, ci scrie heartbeat. Aceste detalii reduc riscul ca platforma să pară funcțională în timp ce componentele critice sunt indisponibile.

\section{Fluxuri concrete de API}
Implementarea API-ului este gândită pentru a păstra rutele apropiate de limbajul produsului. Rutele de learning expun catalog, detaliu de curs, progres, evaluări, înscrieri și administrare de conținut. Rutele de teacher expun workspace-ul profesorului, cererile de înscriere, studenții și candidații la recompense. Rutele de rewards expun istoricul, politicile, candidații, deciziile și fraud blocks. Rutele de wallet expun legarea portofelului, auditul, depunerile, retragerile și taxele. Această separare face API-ul mai ușor de folosit și documentat.

Un avantaj al acestei organizări este că permisiunile pot fi aplicate aproape de resursa protejată. O rută de curs conține identificatorul cursului și poate verifica permisiuni de curs. O rută de organizație conține identificatorul organizației și poate verifica rolul organizațional. O rută de platformă cere drepturi globale. Dacă toate operațiile ar fi grupate generic sub un endpoint de administrare, această claritate s-ar pierde.

\section{Validarea datelor la intrare}
Datele primite prin API trebuie validate înainte de a ajunge în domeniu. Validarea include câmpuri obligatorii, lungimi, enum-uri, identificatori, sume și formate. De exemplu, o sumă de recompensă trebuie să fie pozitivă și compatibilă cu tipul numeric folosit în bază de date. O adresă de portofel trebuie să aibă format acceptabil. O decizie de profesor trebuie să conțină motiv atunci când respinge un candidat. Această validare previne erori mai jos în stivă.

Validarea nu înlocuiește regulile de domeniu. Chiar dacă un request este bine format, operația poate fi invalidă. Un curs poate fi arhivat, o politică poate fi inactivă, un fraud block poate opri recompensa sau actorul poate avea permisiune insuficientă. De aceea, RustLearn separă validarea formei de validarea sensului. Prima aparține stratului HTTP și DTO-urilor, a doua aparține use case-urilor și domeniului.

\section{Implementarea exporturilor și rapoartelor}
Exporturile CSV și rapoartele sunt utile deoarece datele platformei trebuie analizate în afara interfeței curente. Administratorii pot avea nevoie de situații privind recompensele în așteptare, wallet-urile organizațiilor, fraud blocks active sau progresul cursurilor sponsorizate. Implementarea exporturilor trebuie să respecte aceleași permisiuni ca interfața. Un export nu trebuie să devină o scurtătură pentru acces la date nepermise.

Din punct de vedere tehnic, exportul transformă rezultate de use case în format tabular. Această transformare trebuie să fie explicită: coloane stabile, valori clare, date formatate și lipsa secretelor. Pentru o platformă educațională, exportul este important și în procesul de evaluare, deoarece profesorii și organizațiile pot folosi datele pentru decizii. Astfel, raportarea este o funcție de produs, nu doar o utilitate administrativă.

\section{Corelarea frontend-backend}
Frontend-ul și backend-ul trebuie să evolueze împreună. Backend-ul definește contractele de date, iar frontend-ul le transformă în interfețe. Dacă backend-ul returnează stări de recompensă prea tehnice, frontend-ul trebuie să le traducă în mesaje utile. Dacă frontend-ul cere un câmp care nu are sens pentru un rol, contractul API trebuie ajustat sau explicat. Această corelare este vizibilă în tipurile din \texttt{web/src/lib}, unde rutele frontend folosesc funcții dedicate pentru sesiune, teacher, learner, organizații și wallet.

Un beneficiu al acestei structuri este izolarea erorilor. Dacă o rută de reward își schimbă răspunsul, adaptarea se face în biblioteca frontend aferentă, nu în toate componentele. Dacă sesiunea primește o capabilitate nouă, verificarea poate fi reutilizată în mai multe pagini. Această disciplină reduce duplicarea și susține interfețe coerente pentru roluri diferite.

\section{Configurare și separarea mediilor}
Configurarea aplicației este tratată ca parte a implementării, nu ca detaliu secundar. Cheile JWT, conexiunile PostgreSQL, endpoint-urile S3, adresele contractelor, cheia deployer-ului și setările worker-ului trebuie furnizate prin variabile de mediu sau secrete, nu prin cod sursă. Această separare permite rularea aceluiași cod în dezvoltare, testare, staging și producție, cu configurații diferite. În plus, reduce riscul ca secretele să fie comise accidental.

Separarea mediilor este importantă și pentru blockchain. În local se poate folosi Anvil, unde tokenurile nu au valoare reală și contractele pot fi redeployate ușor. În staging se poate folosi o rețea de test, cu date apropiate de producție, dar fără risc financiar major. În producție, contractele trebuie auditate, adresele trebuie controlate și secretele trebuie gestionate prin infrastructură specializată. Această trecere graduală este o condiție pentru o platformă care ar putea administra recompense cu valoare reală.

\section{Calitatea codului și verificări înainte de livrare}
Calitatea implementării este susținută prin format, teste și preflight. Formatul uniform reduce diferențele inutile în review. Testele verifică reguli de domeniu, API-uri, contracte și fluxuri de produs. Preflight-ul combină mai multe verificări pentru a oferi o imagine rapidă asupra stării repository-ului. În plus, migrațiile sunt rulate prin Make/Compose, ceea ce evită dependența de o instalare locală Diesel diferită de mediul proiectului.

Aceste verificări sunt relevante pentru disertație deoarece arată că sistemul nu este doar o descriere teoretică. El are o disciplină de dezvoltare care permite modificări repetate fără pierderea controlului. Într-un proiect cu recompense, permisiuni și contracte, această disciplină este esențială. O eroare de format este minoră, dar lipsa unei culturi de verificare poate ascunde erori de autorizare, migrații incomplete sau stări financiare incorecte.
